% !TEX root = ../main.tex
\subsection{Error Recursions for Refreshing and Mixing}
\label{app:algorithmic-details}

Table~\ref{tab:transport-building-blocks} gives the estimator identities.
Here we record how these operations affect tracking error. Write
$u_{k+1}:=H_{k+1}-\nabla f(x_{k+1})$,
$q_k:=\mathbb E[\|d_{k+1}\|^2]$, and
$v_{k+1}:=\mathbb E[\|u_{k+1}\|^2]$.
Both $u_{k+1}$ and $d_{k+1}$ are conditionally unbiased given
$\mathcal F_k$.

\paragraph{Restarts and probabilistic refreshing.}
Within a transport phase starting at $s$, the martingale identity gives
\begin{equation}
\label{eq:restart-error-window}
E_k=E_s+\sum_{t=s}^{k-1}q_t.
\end{equation}
An exact refresh resets $E_s$ to zero; an independent batch of size
$B$ gives $E_s\leq\sigma^2/B$ under BV. This limits the accumulation
window without undoing the effect of past errors on the iterates.
For PAGE, an independent Bernoulli decision selects a fresh estimate
with probability $p$ and transport otherwise \citep{li2021page}.
Conditioning on that decision yields
\begin{equation}
\label{eq:refresh-error-recursion}
E_{k+1}=(1-p)(E_k+q_k)+p\,v_{k+1}.
\end{equation}
The branches are mutually exclusive, so their noises have no cross
term. MARINA uses the same principle with exact refreshes and
unbiased compressed gradient differences \citep{gorbunov2021marina}.
Unbiased compression ensures the conditional mean, but a similarity
bound also requires control of compression variance and local gradients.

\paragraph{Mixing fresh and transported estimates.}
For MVR/STORM with deterministic fresh weight $\alpha$,
subtracting the true gradient from the update in
Table~\ref{tab:transport-building-blocks} gives
\begin{equation}
\label{eq:mvr-error-recursion}
e_{k+1}=(1-\alpha)e_k+(1-\alpha)d_{k+1}+\alpha u_{k+1}.
\end{equation}
Conditional unbiasedness removes cross terms with $e_k$, hence
\[
\begin{split}
E_{k+1}={}&(1-\alpha)^2(E_k+q_k)+\alpha^2v_{k+1}\\
&+2\alpha(1-\alpha)\mathbb E[\langle d_{k+1},u_{k+1}\rangle].
\end{split}
\]
Unlike probabilistic refreshing, mixing uses both noises together.
With shared sampling their covariance cannot generally be discarded.
For history-dependent weights, this identity holds conditionally;
the weights remain inside expectations when averaging over the history.

\paragraph{Scaled transport within momentum.}
The unclipped MARS estimator in
Table~\ref{tab:transport-building-blocks} interpolates between MSGD
at $\gamma=0$ and MVR at $\gamma=1$. For $0\leq\gamma\leq1$,
\[
g_{k+1}^{\mathrm{MARS}}
=(1-\gamma)g_{k+1}^{\mathrm{MSGD}}
+\gamma g_{k+1}^{\mathrm{MVR}},
\]
where the two candidates use the same incoming estimate, sample,
and fresh weight $\beta$ \citep{venus2026draft}.
Pure transport is the $(\beta,\gamma)=(0,1)$ limit of this expanded
identity, not a substitution into MARS's original rescaled formula,
which is undefined at $\beta=0$.
A scaled difference with $\gamma\ne1$ is generally not an unbiased
estimate of the gradient change. Cached samples, clipping, and
preconditioning also fall outside the pure-transport assumptions.
VENUS changes the coupling of fresh and transported samples, and
therefore their covariance, while preserving shared sampling within
each difference \citep{venus2026draft}.

\subsection{Further Connections to Gradient Transport}
\label{app:related-transport}

\paragraph{Implicit gradient transport.}
IGT \citep{arnold2019reducingvarianceonlineoptimization} averages past
stochastic gradients and compensates for their staleness by evaluating
new gradients at extrapolated points. Unlike TSGD, it does not form
shared-sample differences. Its exact correction is derived for quadratic
components with a common Hessian (their Assumption~3.1), a case where
our $\delta=0$. The quadratic restriction matters: $\delta=0$ alone
also permits a nonlinear common gradient plus an additive random offset,
and does not ensure that extrapolation corrects staleness exactly.
Even in the common-Hessian quadratic case, IGT's correction removes
staleness, not sampling noise. TSGD instead adds sampled gradient
changes to a single estimate, using two pointwise evaluations per
increment rather than one at an extrapolated point.

\paragraph{Hessian-based gradient tracking.}
\citet{gower2018trackinggradientsusinghessian} modify SVRG by using
Hessian-based linear models to update its gradient control variates
away from the reference point. The Hessian tracks changes in gradients
rather than preconditioning the update; diagonal and low-rank
approximations reduce its computational cost. Both this approach and
TSGD account for changing gradients, but through different information:
the former uses a second-order model of reference gradients, whereas
TSGD uses shared-sample differences and recursively transports one
estimate without a Hessian model.

\paragraph{The recursive-gradient estimator used with I-CGM.}
\citet{jiang2026nonconvex} use a recursive-gradient (RG) estimator
within an inexact composite gradient method for federated optimization.
In our notation, their Section~5 recursion is
\[
g_0=\nabla f(x_0),\qquad
g_{k+1}=(1-\beta)g_k+\beta H_k+D_{k+1},
\qquad 0<\beta\leq1,
\]
where $H_k$ is an unbiased estimate at the \emph{previous} point $x_k$
and $D_{k+1}$ is a shared-minibatch gradient difference. As the authors
note, setting $\beta=0$ gives the SARAH recursion, hence pure transport
with exact initialization. This identifies the estimator, not the outer
I-CGM iterate update, with TSGD. Their average dissimilarity assumption
is centered similarity under uniform finite-sum sampling.
For effective minibatch similarity constant $\delta_m$, Lemma~5.2
weights the sum of squared displacements by
$2\delta_m^2/(2\beta-\beta^2)$. When $\delta_m>0$, this coefficient
diverges as $\beta\downarrow0$, so that lemma does not cover pure
transport. The finite-horizon analysis in Section~\ref{sec:tsgd-convergence}
instead controls accumulation without requiring a positive mixing weight.
